\section{Limitations}
\label{sec:limitations}

\textbf{One simulator, small pools.} All evidence comes from our generator. Outside the
noisy-lead stratum its lead times are deterministic, so a lead-time shift appears in the receipts
without noise, and that is where the gain sits: at $p/h = 4$ lead-time shifts, a sixth of the shocked episodes, carry
about 171 of the 178 per-episode gain over arm~1 with Gemini and 157 of 163 with Grok. With
noisy lead times the gain over arm~1 was not confirmed on eight units, and the one-period
shifts lost on average (descriptive). The method has not been run on InventoryBench's own instances, a third of
which have stochastic lead times with lost shipments; the fresh pools have 48 units (the confirmation
registered 72) because the alert bank yields distinct texts for only eight units per family.

\textbf{What the theorems cover.} Theorems~\ref{thm:exposure}
and~\ref{thm:regret} assume an exact null, and the arm as run lies outside it: its
demand null is estimated from the pre-proposal prefix, and its arrival null is a
diffuse delay law while our delays are deterministic outside the noisy-lead stratum.
A null Monte Carlo of the frozen arm (released with our code) finds no cell more
than two standard errors above a bound where the premise holds. With the plug-in
null and a first firing at period 10, in the cells matched to our null twins, it
measures $\Expect_0[\Pi_\sigma]$ at 2.0 to 2.7 times $A$ and
$\Prob_0(\sup_t \Pi_t \ge 1/2)$ at 2.3 to 3.7 times its bound (up to 3.2 and 4.7
in other cells); the arrival mismatch alone lifts exposure at a fixed decision to
at most 1.24 times its budget. The regret theorem bounds a one-period newsvendor
surrogate only.

\textbf{What the content-free control measures.} The control calls no model but
is timed by the same alerts, so it removes the text, not the alert channel. Its
null commitment also overstates action: a pulse hypothesis keeps its \eprocess{}
after its window closes, so its mass can linger on a null episode after it stops
moving the order. Post hoc, crediting each period's mass to its leading hypothesis,
pulse-led periods carry 0.69 of the control's 0.83 and 0.32 and 0.38 of the hedge's 0.52
and 0.46, so most of what the model's label removes is that mass.

\textbf{Status of the evidence.} Relative to Gate~3 this study is exploratory, and
Gate~3's decision, \emph{kill-or-reframe}, stands. The method was developed on the spent
pilot episodes; both later stages were registered, code frozen by hash, before their live
runs. The plan discloses a structural test on six first-stage units, an interim look at one
run, three departures from the original registration (net reward as endpoint, action before
any activation threshold, fresh units outside the frozen manifest), an audit that reproduced
every registered first-stage number, and the second stage's pre-run amendments, which changed
no hypothesis or cap (one clipped the noisy stratum's lead times). The second stage did not
establish null safety at $p/h = 19$, nor non-inferiority against arm~1 at $p/h = 1$ with Grok.
